\documentclass[11pt]{article}

\usepackage[T1]{fontenc}
\usepackage[utf8]{inputenc}
\usepackage{newtxtext,newtxmath}
\usepackage[margin=0.85in]{geometry}
\usepackage{xcolor}
\usepackage{setspace}
\usepackage{parskip}
\usepackage[colorlinks,urlcolor=blue,linkcolor=black]{hyperref}

\setstretch{1.08}
\definecolor{responseblue}{RGB}{0,77,128}
\newcommand{\response}{\par\noindent\textbf{Response:} }
\newcommand{\changes}{\par\noindent\textbf{Changes in the revised manuscript:} }

\title{Response to Reviewers\\[0.4em]
\large COMMAT-D-26-03061}
\author{Yuhan Sun and Xiao Huang}
\date{24 August 2026}

\begin{document}
\maketitle

\noindent\textbf{Revised title:} Local-environment-dependent transferability of
a fine-tuned MACE potential for lithium adsorption on graphene defects

\section*{General response}

We thank the Associate Editor and both reviewers for their constructive
assessment. We agree with the central criticism that the submitted manuscript
read too much like a technical audit and did not lead with a scientific
question or physical findings. We therefore performed a substantive
reorganization rather than a local language edit.

The revised manuscript asks how graphene defects and a local Si$_4$ motif alter
dilute Li adsorption and how those chemically distinct environments affect the
transferability of a fine-tuned foundation MLIP. The revised paper now leads
with PBE-D3(BJ) adsorption-energy anchors and direct DFT force tests on displaced
configurations. Dataset-split audits, $E_0$ details, fixed-path scans, trajectory
traces, and snapshot traceability are retained in Methods or Supporting
Information, where they support the science without becoming the narrative.

We also added an explicit comparison with DeePMD-kit/DP-GEN, clarified that the
three-model committee is a qualitative seed-sensitivity diagnostic rather than
a calibrated uncertainty ensemble, and stated that the DFT-validated snapshots
have not yet entered a closed retraining cycle. The Methods section now records
the inherited MACE architecture, complete fine-tuning settings, checkpoint
selection, deterministic $E_0$ solve, dataset scope, group-aware split, and DFT
snapshot protocol. The manuscript has continuous line numbering for review.

\section*{Reviewer 2}

\subsection*{Comment 1: Relationship to concurrent-learning and active-learning workflows}

\begin{quote}
The relationship between the proposed validation-first workflow and established
concurrent-learning or active-learning MLIP workflows, particularly DeePMD-kit
and DP-GEN, should be discussed. The common elements and the distinction from a
closed iterative learning cycle should be stated explicitly; a comparison table
may be useful.
\end{quote}

\response We agree. The revised Introduction now states that the
explore--label--retrain logic is established rather than novel. It describes the
DP-GEN sequence of exploration, ensemble model-deviation screening,
first-principles labeling, and retraining, and cites both DeePMD-kit and DP-GEN.
We then state the precise difference: our study does not use a quantitative
model-deviation threshold and does not complete a subsequent retraining cycle.
Its function is to gate scientific claims and identify targeted future labels.

\changes Introduction, p.~2, lines 46--54; new references 14 and 15, p.~10,
lines 315--319. A six-row comparison of model, exploration, selection signal,
DFT step, retraining loop, and purpose is provided in Supporting Information
Table~2 (p.~2).

\subsection*{Comment 2: Precision of the novelty statement}

\begin{quote}
The novelty should not be presented as the general use of ML-driven exploration
followed by first-principles labeling. The contribution should be defined more
precisely in terms of validation of downstream scientific claims and narrowing
claims when validation fails.
\end{quote}

\response We agree and have removed ``validation-first'' from the title. The
revised manuscript no longer claims a new active-learning architecture. Its
scientific contribution is now the environment-resolved finding that defect
chemistry changes both Li adsorption and MLIP reliability. The methodological
point is stated narrowly: independent DFT tests determine which adsorption and
transferability claims are supported.

\changes New title and Abstract, p.~1, lines 1--27; revised novelty statement,
p.~2, lines 46--60; scientific question and result, pp.~2--3, lines 61--69.

\subsection*{Comment 3: Role of the three-member MACE committee}

\begin{quote}
Please clarify whether committee disagreement quantitatively selected
configurations for DFT or served only as a qualitative model-sensitivity
diagnostic. If no force-deviation criterion comparable to concurrent-learning
thresholds was used, this should be stated.
\end{quote}

\response The committee was used only to evaluate sensitivity to independent
fine-tuning seeds. No calibrated uncertainty estimate, committee-disagreement
threshold, or force-deviation threshold was used to select DFT configurations.
We now state this directly and avoid calling the committee an uncertainty
ensemble.

\changes Potential development and validation workflow, pp.~5--6, lines
154--162. Seed-resolved metrics remain in Supporting Information Section~1.

\subsection*{Comment 4: Limitations of the grouped split}

\begin{quote}
The group-aware split is appropriate, but the five-configuration test set is a
leakage audit rather than a complete transferability benchmark. Avoid wording
such as ``improved transferability'' unless supported by additional independent
configurations.
\end{quote}

\response We agree. The five grouped test configurations are never used as a
population-level transferability estimate. The revised text calls this split a
leakage audit and uses the original same-workflow metric only as an interpolation
baseline. Transferability is assessed instead by direct DFT force comparisons
for displaced configurations.

\changes Methods, p.~5, lines 124--128; Results, p.~7, lines 199--213. The
Abstract explicitly calls the 285.2-to-20.1 meV~\AA$^{-1}$ change an
``interpolation improvement but not independent transferability'' (p.~1, lines
19--21).

\subsection*{Comment 5: Path spans are not migration barriers}

\begin{quote}
Keep the distinction between fixed-geometry path spans and migration barriers
consistent throughout the manuscript, figures, Supporting Information, and data
files. Do not use barrier, activation-energy, or migration-energy terminology
for these descriptors.
\end{quote}

\response We agree. The main text reports only fixed-geometry path-scan
descriptors and endpoint energy differences and makes no kinetic claim. The
complete scans are moved to the Supporting Information. We also changed the
curated public CSV columns from the historical \texttt{barrier\_from\_*} names
to \texttt{max\_minus\_start\_ev} and \texttt{path\_span\_ev}.

\changes Methods, p.~4, lines 108--117; Scope and limitations, p.~8, lines
238--247; Supporting Information Figure~3 and Table~3 (p.~4); curated file
\path{submission_data/results/fixed_path_descriptors.csv}.

\subsection*{Comment 6: Interpretation of molecular-dynamics results}

\begin{quote}
The conservative treatment of the seed-dependent MSD traces is appropriate.
Please state explicitly that successful completion of an ML-driven MD run does
not establish physical reliability, especially because the direct DFT force
errors are much larger than the same-workflow test errors.
\end{quote}

\response We added the requested statement. Trajectories are used only to
generate displaced configurations; completion is described as numerical
execution rather than physical validation. The manuscript does not report a
diffusion coefficient.

\changes Finite-temperature snapshot generation, p.~6, lines 163--178,
especially lines 167--169. The trajectory traces and their seed dependence are
now in Supporting Information Figure~4 and Table~4 (p.~5).

\subsection*{Comment 7: Emphasize direct DFT validation}

\begin{quote}
The direct DFT validation of high-displacement configurations is among the
strongest results and should be emphasized in the Discussion and Conclusion.
Keep the interpretation tied to demonstrated DFT--MACE discrepancies rather
than displacement magnitude alone.
\end{quote}

\response We agree. The direct DFT comparison is now the central panel of
Figure~2 and the central Results subsection. The interpretation is tied to the
measured force errors and to a specific structural regime: short reconstructed
C--C contacts in the selected monovacancy frames. We do not infer that all
high-displacement configurations are out of domain.

\changes Abstract, p.~1, lines 19--27; Figure~2 and interpolation boundary,
p.~7, lines 199--213; environment-resolved DFT result, p.~8, lines 214--237;
Conclusion, p.~9, lines 248--260.

\subsection*{Comment 8: Whether DFT snapshots entered another training cycle}

\begin{quote}
Please clarify whether the DFT-validated MD configurations were incorporated
into a subsequent training cycle. If they are future targets, state that a
closed active-learning/retraining cycle is beyond the scope of the study.
\end{quote}

\response They were not incorporated into a subsequent MACE training cycle. We
now state this in both Methods and Results. The DFT labels identify the next
chemistry-specific training targets, but the current study does not claim a
closed active-learning loop.

\changes Methods, p.~6, lines 173--178; Results, p.~8, lines 229--237.

\section*{Reviewer 3}

\subsection*{Comment 1: Rewrite the Abstract as a scientific summary}

\begin{quote}
The Abstract reads like a technical audit log. It should follow a conventional
background--innovation--validated results--implications structure, emphasize the
main physical findings, and move pre-audit metrics and detailed counts to the
main text.
\end{quote}

\response We agree and rewrote the Abstract completely. It now begins with the
scientific problem, defines the five local environments and DFT--MACE approach,
reports defect-dependent Li adsorption and environment-resolved DFT force
errors, and closes with the implication for targeted data expansion. It does
not list job completion, excluded NEB claims, checkpoint history, or snapshot
counts.

\changes Revised Abstract, p.~1, lines 9--27. Detailed split, $E_0$, path,
trajectory, and traceability material is now in Methods or Supporting
Information.

\subsection*{Comment 2: Figure quality, contrast, and panel labels}

\begin{quote}
Figure 5 has poor curve contrast and indistinguishable line styles. Figure 4
lacks subpanel labels. Use colorblind-friendly palettes, distinct markers or
line types, consistent fonts, and at least 300 dpi for raster graphics.
\end{quote}

\response We reorganized the figures rather than simply shrinking the old
multi-panel diagnostics. The main manuscript now contains only two scientific
figures: a high-resolution structure overview and a vector summary of Li
adsorption and environment-resolved DFT force errors. The former fixed-path
figure is moved to Supporting Information Figure~3, where all ten panels are
labeled (a)--(j). The MSD traces are moved to Supporting Information Figure~4
and redrawn with a colorblind-friendly palette, independent solid/dashed/dotted
line styles, three marker shapes, explicit (a)/(b) panel labels, and a vector
PDF; the companion PNG is 3150$\times$1260 pixels at 300 dpi.

\changes Main Figures~1--2; Supporting Information Figures~2--4 (pp.~3--5),
especially the revised path panels on p.~4 and MSD panels on p.~5.

\subsection*{Comment 3: Language and academic tone}

\begin{quote}
The manuscript requires extensive language polishing. Informal or internal-note
phrasing and repeated First/Second/Third/Finally constructions should be
removed, and the entire paper should use precise scientific English.
\end{quote}

\response We agree. The title, Abstract, Introduction, Results, and Conclusion
were rewritten around a continuous scientific argument. Internal process
phrases such as ``language is restored'' and ``runtime-completion stress tests''
have been removed. Operational logs, status language, and repeated lists of
failed gates no longer appear in the scientific narrative. Limitations are
consolidated in one subsection rather than repeated as debug commentary.

\changes Title and Abstract, p.~1, lines 1--27; Introduction, pp.~2--3, lines
30--69; Results and Discussion, pp.~6--8, lines 179--247; Conclusion, p.~9,
lines 248--260.

\subsection*{Comment 4: Complete the potential-development methodology}

\begin{quote}
The manuscript should give the MACE hyperparameters and training protocol, the
criterion for the foundation-assisted $E_0$ estimate, justification for the
273-frame dataset, and an explicit group-held-out split strategy. A dedicated
Potential Development Workflow subsection or flowchart is recommended.
\end{quote}

\response We added a dedicated ``Potential development and validation
workflow'' subsection. It records the inherited architecture, cutoff,
interaction layers, irreducible representations, correlation order, angular
limit, precision, batch size, epochs, learning rate, optimizer, weight decay,
EMA, SWA transition, stage weights, trainable components, selected checkpoints,
and all three $E_0$ offsets. We clarify that the architecture was inherited
unchanged rather than screened against the test set. The $E_0$ estimate is a
deterministic full-rank linear solve, not an iterative optimization; acceptance
requires finite coefficients and elemental rank 3/3.

The dataset is now described as a deliberately compact sample of five local
chemical families and selected displacement coordinates, not exhaustive
configurational coverage. The exact group-aware split algorithm and the
263/5/5 outcome are stated. The Supporting Information adds the complete split
metrics, seed sensitivity, $E_0$ details, and a workflow-comparison table.

\changes Dataset scope, p.~4, lines 108--114; group-aware split, pp.~4--5,
lines 118--128; architecture and training protocol, p.~5, lines 129--162;
snapshot validation protocol, p.~6, lines 163--178; Supporting Information
Tables~1--2 (p.~2).

\section*{Closing statement}

We believe these revisions address every point raised by both reviewers and
convert the paper from a workflow audit into a focused scientific study of
defect-dependent Li adsorption and local-environment-dependent MLIP
transferability. All numerical claims in the revised manuscript were checked
against the archived evidence tables, and the source, figures, Supporting
Information, and compact reproducibility package are supplied with the
resubmission.

\end{document}
